\chapter{Solución Propuesta}
\label{ch:solucion}

Este proyecto propone el desarrollo de una metodología de emulación física y caracterización para 
superar las limitaciones del hardware cerrado y disminuir los riesgos de la fabricación en 
silicio sin el previo modelado del sistema necesario. Para esto, se proponen cuatro etapas 
para la solución.

\section{Primera Etapa}

La primera etapa de la solución propuesta para el proyecto se enfoca en la selección de 
una arquitectura de lógica reconfigurable de código abierto parametrizable, y su 
implementación mediante un banco de pruebas para comprobar su funcionamiento correcto. De 
esta manera se elimina la dependencia del flujo de herramientas de un fabricante específico 
y se asegura la portabilidad del proyecto a diferentes plataformas, al igual que se garantiza 
la continuidad del proyecto en el caso de querer profundizar en el futuro.

\section{Segunda Etapa}

La segunda parte de la solución propuesta para el proyecto se centra en la instanciación de 
la arquitectura de lógica reconfigurable seleccionada en la primera etapa sobre una FPGA 
comercial. Mediante esta etapa se busca un medio para extraer en próximas etapas las métricas 
de rendimiento y uso de recursos como el área equivalente, el consumo de potencia y los 
retardos de las rutas críticas.

\section{Tercera Etapa}

La tercera fase de la solución propuesta para el proyecto busca implementar dos circuitos 
digitales que estresen la arquitectura de lógica reconfigurable seleccionada en la primera 
etapa de maneras diferentes. Para esto se propone la implementación de un filtro digital 
de orden bajo y un multiplicador de enteros. De esta manera se busca obtener métricas de 
rendimiento para su comparación y se garantiza la funcionalidad del flujo de síntesis y 
ruteo de la arquitectura de lógica reconfigurable seleccionada.

\section{Cuarta Etapa}

La cuarta etapa de la solución propuesta para el proyecto tiene como objetivo la consolidación 
de los datos obtenidos en la tercera etapa observando cómo las métricas de rendimiento se 
compensan entre sí, para determinar si la arquitectura de lógica reconfigurable seleccionada 
es adecuada para su uso en aplicaciones de alto rendimiento o si es necesario realizar 
modificaciones según la aplicación en la que se utilice. Esta etapa se enfoca en la 
interpretación de los datos y cómo su análisis puede guiar el diseño de futuras arquitecturas 
de lógica reconfigurable para su impresión en silicio.